Nothing in the empirical sections depends on this one. It is here because it
is how we came to look at $\bar H$, and because a reader is entitled to know
that the route was a theory rather than a hunch.

We treat merging as lossy compression. A rate-$R$ merging code is an encoder
and decoder pair that between them turn $T$ task vectors into one weight
vector using $R$ bits, with both ends allowed to see $\boldsymbol H$ for free,
and $\mathcal{D}^\star(T,d,r,B,R)$ is the best worst-task distortion any such
code achieves against the hardest admissible task tuple. The formal
definitions are in Appendix~\ref{app:rd}.

Against a rotation-invariant hard distribution, every merging code obeys
\begin{equation}\label{eq:bound-brief}
  \mathcal{D}^\star
  \;\geq\;
  \underbrace{B^2\Bigl(1 - \frac{\E[\deff]}{Tr}\Bigr)}_{\text{irreducible floor}}
  \;+\;
  \underbrace{c_{\mathrm{TQ}}\cdot\frac{B^2\E[\deff]}{Tr}\,
  2^{-2R/\E[\deff]}}_{\text{rate term}},
\end{equation}
and in the regime where the floor vanishes an explicit construction matches
the exponent up to a universal constant. Statements, hypotheses and proofs are
in Appendix~\ref{app:rd} (Theorems~\ref{thm:general} and~\ref{thm:achv}); the
proofs themselves are in Appendix~\ref{app:proofs}.

\paragraph{What it did.} It put $\deff$, and therefore the geometry of the
union of the task subspaces, at the centre of the problem. That is what sent
us to measure principal cosines between task row spaces, which is the
measurement the rest of the paper is built on and the one we expect a
practitioner to use. Nothing else we tried pointed there.

\paragraph{What it does not do, stated here rather than in a limitations
list.} It derives none of the empirical results that follow, and three
specific things are wrong with using it as support.

Its floor term is \textbf{exactly zero} on all four base models under every
initialisation scheme we tried, because $T$ rank-$r$ adapters with $Tr \leq d$
are linearly independent even when their subspaces are nearly collinear. Near
collinear is not collinear. So on real LoRA cohorts the floor carries no
operational content at all (\S\ref{subsec:floor-zero}). The trap worth
naming is that substituting a soft participation-ratio surrogate for the
rank makes the floor look large and cohort-dependent when it is neither.

Its rate term rests on Assumption~\ref{asm:epr}, which we state and do not
prove, and on a constant we import from TurboQuant~\citep{zandieh2025turboquant}
rather than establish. Our attempt to measure the predicted exponent has too
little dynamic range to do so, so we cannot claim the term is even consistent
with our data (Appendix~\ref{app:rate}).

And Theorem~\ref{thm:general} constrains a $w^\star$ carrying no rank
constraint, while every merge we run is truncated to rank $r$. At that step it
bears on our measurements by analogy and not by derivation
(\S\ref{subsec:conditioning}).

We think the lens earned its keep. But it is a lens, and everything in
\S\ref{sec:regimes} and \S\ref{sec:methods} stands or falls on its own
measurements.
